\subsection{Assignment 9 -- Analisi Summer - Winter Score per regione}

Per affrontare l’Assignment 9 è stato progettato un flusso di dati in
\textit{SQL Server Integration Services} (SSIS) a partire dalla fact table dei
brani pubblicati. In una prima fase, il flusso è stato filtrato tramite una
trasformazione di \textit{Conditional Split} per mantenere esclusivamente le
righe in cui l’artista ricopre il ruolo di \textit{main artist}, garantendo
coerenza con la granularità dell’analisi.

A partire da queste righe, sono state introdotte tre trasformazioni di
\textit{Lookup}. Il primo Lookup con la dimensione \textit{Artist} consente di
recuperare l’attributo \textit{GeoKey}, necessario per il secondo Lookup con la
dimensione \textit{ArtistGeography}, che permette di associare a ciascun artista
la regione e la nazione di riferimento. L’ultimo Lookup è stato effettuato sulla
dimensione \textit{Date} al fine di ricavare la stagione di pubblicazione del
brano. Al termine di questa fase, ogni riga del flusso risulta caratterizzata
congiuntamente dalla regione dell’artista e dalla stagione di pubblicazione.

Il flusso viene quindi suddiviso in due percorsi distinti, corrispondenti
rispettivamente ai brani pubblicati in estate e a quelli pubblicati in inverno.
Ciascun percorso è aggregato separatamente per regione tramite una
trasformazione di \textit{Aggregate}, calcolando da un lato il numero di brani
estivi e la somma degli \textit{Streams1Month}, e dall’altro le corrispondenti
misure per i brani invernali.

Poiché la trasformazione di \textit{Merge Join} richiede flussi ordinati, entrambi
i percorsi sono sottoposti a un’operazione di \textit{Sort} prima di essere
ricombinati. Nel \textit{Merge Join} è stata utilizzata una \textit{full outer
join}, così da preservare anche le regioni presenti esclusivamente in uno dei
due flussi stagionali.

L’output risultante contiene, per ciascuna regione, quattro misure: il numero di
brani estivi, il numero di brani invernali, la somma degli stream estivi e la
somma degli stream invernali. A partire da questi valori, è stata introdotta una
trasformazione di \textit{Derived Column} per calcolare i due indicatori
richiesti: il \textit{summer--winter score} basato sul numero di brani e quello
basato sul numero di stream.

Entrambi i calcoli includono una gestione esplicita dei casi in cui il
denominatore risulta pari a zero, restituendo un valore \textit{NULL} al fine di
evitare divisioni invalide. Questa scelta previene la produzione di misure
distorte e rende gli indicatori più robusti e interpretabili dal punto di vista
analitico.

Il flusso così completato viene infine indirizzato verso una \textit{Flat File
Destination}, che rappresenta l’output finale. 
